\section{Performance Evaluation}
\label{sec:exp}

6FLAS is evaluated against four baselines using a reproducible two-tier setup.

\subsection{Evaluation Methodology}
\label{sec:setup}

An \emph{analytical processing model} charges each mechanism the cost of its
per-hop primitives in CPU cycles on a 1\,GHz space-grade core ($1$ cycle
$=1$\,ns), using primitive costs drawn from standard embedded software-crypto
microbenchmarks; parsing, MAC, encryption, replay, and association-lookup terms
are listed explicitly so the model can be audited and re-parameterized.

An \emph{event-driven latency model} routes random
source--destination pairs over a $+$Grid topology and accumulates free-space
propagation, per-hop security processing, and light $M/M/1$ queueing.
The constellation is a Starlink-like shell: $h=550$\,km, $P=72$ planes,
$S=22$ satellites per plane, $N=1584$ nodes, intra-plane ISL $\approx1970$\,km
and inter-plane ISL $\approx604$\,km~\cite{del2019,bhattacherjee2022}. The entire evaluation is regenerated
deterministically (seed fixed) by a single script released with the paper.

A C microbenchmark (\texttt{proto/bench\_flas.c}, co-released) provides
cycle-accurate empirical measurements on commodity x86. RDTSC-bracketed loops
of $10^6$ iterations measure SipHash-2-4 (FLAS), HMAC-SHA256 (SRv6-Sec), and
HMAC-SHA256 + software AES-128-CBC (IPsec-like) over the same covered spans as
the analytical model. Because the benchmark runs on commodity x86 rather than
space-grade hardware, absolute cycle counts
are platform-dependent. The microbenchmark's primary contribution is confirming the structural
ordering $\text{FLAS} \ll \text{SRv6-Sec} \ll \text{IPsec}$, which is
architecture-independent.

Four baselines are used for comparison: \textbf{Plain SRv6} (routing only, no
security; a lower bound on processing), \textbf{IPsec ESP} (transport mode,
HMAC-SHA256 + AES-128-CBC, per-flow SA), \textbf{DTLS 1.2} (record layer, same
primitives), and \textbf{SRv6-Sec (HMAC)}, which verifies an HMAC-SHA256 over
the routing header per hop in the spirit of the SRH HMAC TLV~\cite{rfc8754}.
Four categories of metrics are reported: \textbf{per-hop cost} (processing time,
communication overhead, on-board state); \textbf{scalability} (forwarding
throughput vs.\ concurrent-flow count); \textbf{dynamics} (end-to-end latency
and control-plane handover cost); and \textbf{sensitivity} (payload-size
independence, tag-width tradeoff, and energy).

\cref{tab:cyc} lists the per-primitive cycle costs used by the processing model,
normalized to a $1$\,GHz core ($1$ cycle $=1$\,ns). Costs follow widely reported
embedded software-crypto ratios: symmetric MAC/cipher work is charged per
processed byte or per $16$-byte block, with fixed setup terms for HMAC and AES.
The authenticated routing context for 6FLAS includes the IPv6 source and destination addresses ($32$\,B), Flow Label, the active SRv6 SID, and the 6FLAS
metadata block, totaling a covered span processed by the MAC. The SRv6 HMAC baseline covers a similar span. IPsec/DTLS additionally process a
$512$\,B payload, reflecting their payload encryption. All constants and the
resulting figures are produced by the released script.
Absolute values are analytically projected and depend on the chosen primitive costs;
the structural ordering of mechanisms and FLAS's payload-independence are robust
to reasonable re-parameterisation.

\begin{table}[t]
\centering
\caption{Per-primitive cost constants (cycles) used in the processing model.}
\label{tab:cyc}
\renewcommand{\arraystretch}{1.15}
\begin{tabular}{lr}
\toprule
\textbf{Primitive} & \textbf{Cost (cycles)} \\
\midrule
Parse IPv6 base header        & 40 \\
Parse SRH + one SID           & 60 \\
Flow-Label window select      & 25 \\
Replay sliding-window test    & 30 \\
SipHash keyed MAC (per byte)  & 2.1 \\
HMAC-SHA256 setup             & 900 \\
HMAC-SHA256 (per byte)        & 11.0 \\
AES-128-CBC (per 16\,B block) & 220 \\
AES setup                     & 350 \\
IPsec SA / SPD lookup         & 180 \\
DTLS record bookkeeping       & 1500 \\
\bottomrule
\end{tabular}
\end{table}

\subsection{Data-Plane Overhead}
\label{sec:results}

\cref{fig:proc} shows per-hop processing time. 6FLAS verifies a packet in about
$264$\,ns, dominated by the short keyed MAC over the authenticated routing context (a fixed set of immutable fields: 32B addresses, 16B SID, and metadata) and
the $O(1)$ replay-window probe. IPsec ESP costs about $14.6\,\mu$s and DTLS
about $15.6\,\mu$s, since both run a full HMAC-SHA256 plus AES-CBC over the
payload and pay an SA/record lookup; the SRv6 HMAC baseline costs about
$1.59\,\mu$s because of its full SHA-256 digest. 6FLAS is thus $55\times$
faster than IPsec ESP and $6.0\times$ faster than the SRv6 HMAC baseline.
The C microbenchmark measures FLAS at $152$ cycles/packet, SRv6-Sec
at $4{,}091$ cycles, and IPsec-like
at $209{,}954$ cycles. The FLAS:SRv6-Sec ratio on commodity x86
($26.9\times$) exceeds the analytical $6.0\times$ because SipHash-2-4
is significantly faster per byte than HMAC-SHA256, amplifying the ratio beyond
what the conservative space-grade model charges. The structural ordering
$\text{FLAS} \ll \text{SRv6-Sec} \ll \text{IPsec}$ is confirmed.

\begin{figure}[t]
\centering
\includegraphics[width=0.92\columnwidth]{figures/proc_overhead.pdf}
\caption{Per-hop security processing time on a 1\,GHz space-grade core.}
\label{fig:proc}
\end{figure}

In communication overhead (\cref{fig:header}), 6FLAS adds only $8$ bytes (the
metadata block) and reuses the existing Flow Label field at no extra cost, whereas
IPsec ESP adds $38$ bytes (ESP header, IV, ICV, trailer), DTLS adds $37$, and the
SRv6 HMAC TLV adds $36$ (4-byte TLV header + 32-byte digest). 6FLAS therefore
cuts header overhead by about $76\%$ relative to IPsec ESP, which matters on the
bandwidth-limited ISL.

\begin{figure}[t]
\centering
\includegraphics[width=\columnwidth]{figures/header_overhead.pdf}
\caption{Per-packet byte overhead added by each mechanism on top of IPv6+SRH.}
\label{fig:header}
\end{figure}

\subsection{Scalability Under Dynamic Topology}
\label{sec:dyn}

\cref{tab:state} lists per-flow on-board state. 6FLAS carries \emph{zero}
per-flow state: the only mutable on-board storage is the per-$(\textsf{domain}, \textsf{direction})$ replay window. For $D$ domains and 6 directions, the total state is
$D \times 6 \times 32$\,B. For $D=256$ domains, this is $48$\,KB, independent of $|\setF|$. SRv6-Sec (HMAC) requires $96$\,B per
flow: a 32-byte per-flow HMAC key, a 32-byte cached inner-pad
(H(K $\oplus$ ipad)), and a 32-byte sliding-window counter block, so that the
per-hop HMAC can avoid full key re-derivation on every packet. This compares
with IPsec's $256$\,B SA (key material, selectors, replay window, SA
identifiers) and DTLS's $320$\,B session record per flow.
\cref{fig:scal} sweeps the concurrent-flow count from $10$ to $10^6$ and plots
single-core forwarding throughput. Stateful baselines degrade as the working set
grows (association-lookup cache pressure), whereas 6FLAS and Plain SRv6 remain flat because their state is independent of flow count.
At $10^6$ flows 6FLAS sustains about $3.8$\,Mpps on one core, analytically ${\approx}55\times$
the throughput of IPsec ESP.

\begin{table}[t]
\centering
\caption{Per-flow on-board security state.}
\label{tab:state}
\renewcommand{\arraystretch}{1.2}
\begin{tabular}{lcc}
\toprule
\textbf{Mechanism} & \textbf{State / flow (B)} & \textbf{Reduction vs.\ IPsec} \\
\midrule
IPsec ESP        & 256 & --- \\
DTLS 1.2         & 320 & $-25\%$ (worse) \\
SRv6-Sec (HMAC)  & 96  & $63\%$ \\
\textbf{6FLAS}    & \textbf{0\,{\scriptsize(48\,KB total for 256 domains, domain-scoped)}} & $\mathbf{100\%}$ \\
\bottomrule
\end{tabular}
\end{table}

\begin{figure}[t]
\centering
\includegraphics[width=0.95\columnwidth]{figures/scalability.pdf}
\caption{Single-core forwarding throughput versus number of concurrent flows.}
\label{fig:scal}
\end{figure}

\cref{fig:latcdf} shows the latency CDF over the constellation for 6FLAS, IPsec
ESP, and Plain SRv6. End-to-end latency is largely determined by free-space propagation and queueing:
6FLAS's median ($\approx99$\,ms) and
$95$th percentile ($\approx189$\,ms) are statistically indistinguishable from
Plain SRv6 ($\approx98$\,ms median). IPsec ESP's larger per-hop processing
accumulates across hops and adds a visible tail at the high percentiles.

\begin{figure}[t]
\centering
\includegraphics[width=0.95\columnwidth]{figures/latency_cdf.pdf}
\caption{End-to-end latency CDF over the 1584-satellite constellation.}
\label{fig:latcdf}
\end{figure}

\medskip
The LEO data plane is dominated by churn: orbital motion continually
makes and breaks ISLs, so flows are rerouted onto new satellites many times per
orbit. A stateful mechanism must migrate or re-establish its per-flow security
association on the new node, which means a fresh handshake (an asymmetric
operation) and an SA/session install; we charge IPsec a $\sim$IKEv2-class
handshake and DTLS a comparable record-layer handshake, and the SRv6 HMAC
baseline a per-flow key re-derivation. Although techniques such as session resumption or pre-provisioned SAs can reduce handshake cost, the state-migration overhead and SA fanout management remain non-trivial at million-flow scale, especially when topology changes are frequent and predictable lead times are short. 6FLAS migrates \emph{nothing}: the
security context travels in the packet, so the next satellite verifies the very
next packet with the domain key it already holds.

\cref{fig:handover} sweeps the number of flows that must be rerouted within a
one-second control window and reports the aggregate control-plane CPU time this
imposes on the new node. At $10^4$ rerouted flows, IPsec must spend
$35$\,s of control-plane CPU time and DTLS about $38$\,s,
the SRv6 HMAC baseline about
$1.2$\,s, while 6FLAS spends none because no
per-flow security context exists on the satellite.

\begin{figure}[t]
\centering
\includegraphics[width=0.95\columnwidth]{figures/handover_cost.pdf}
\caption{Control-plane CPU time to re-secure flows after an ISL handover versus rerouted-flow count.}
\label{fig:handover}
\end{figure}

\subsection{Sensitivity Analysis}
\label{sec:sensitivity}

6FLAS authenticates only the authenticated routing context
(\cref{tab:cyc}) — a fixed set of immutable routing fields — so its per-hop cost is constant regardless of how large the
data packet is. IPsec and DTLS authenticate and encrypt the payload, so their
per-hop cost grows linearly with packet size. \cref{fig:payload} sweeps the
payload from $64$ to $1500$\,bytes: 6FLAS and the SRv6 HMAC baseline are flat
($264$\,ns and $1.59\,\mu$s), whereas IPsec rises from about $4.6\,\mu$s to
$39\,\mu$s and DTLS similarly. The
gap widens with
packet size, reaching $148\times$ at $1500$\,bytes. Authenticating
routing metadata rather than payload is what makes 6FLAS's cost insensitive to
the traffic it carries.

\begin{figure}[t]
\centering
\includegraphics[width=0.95\columnwidth]{figures/payload_sensitivity.pdf}
\caption{Per-hop processing time versus payload size.}
\label{fig:payload}
\end{figure}

The truncated tag width $t$ sets both the single-packet forgery bound
($2^{-t}$, \cref{prop:forgery}) and the metadata size. \cref{tab:tag} tabulates
the tradeoff. Widening the tag from the $32$-bit default to $64$ bits tightens
the forgery bound from $2^{-32}$ to $2^{-64}$ at a cost of only $4$ extra bytes
per packet and $8$\,ns of extra processing, because the keyed MAC already runs
over the authenticated routing context and only the tag-comparison span grows. A deployment can
thus pick its security margin at a cost of only $4$ extra bytes and ${\approx}8$\,ns per packet; $32$ bits are used as a balance
between the bound and header compactness, and note that even the $64$-bit
variant still adds only $13$ bytes, about a third of IPsec's $38$.

\begin{table}[t]
\centering
\caption{Tag-width tradeoff: forgery bound, metadata size, and per-hop cost.}
\label{tab:tag}
\renewcommand{\arraystretch}{1.2}
\begin{tabular}{cccc}
\toprule
\textbf{Tag (bits)} & \textbf{Forgery bound} & \textbf{Metadata (B)} & \textbf{Proc.\ (ns)} \\
\midrule
16 & $2^{-16}$ & 7  & 258 \\
24 & $2^{-24}$ & 8  & 262 \\
\textbf{32} & $\mathbf{2^{-32}}$ & \textbf{9}  & \textbf{264} \\
48 & $2^{-48}$ & 11 & 269 \\
64 & $2^{-64}$ & 13 & 273 \\
\bottomrule
\end{tabular}
\end{table}

Energy is a first-order constraint on solar/battery-limited satellites. Space-grade radiation-hardened processors (e.g., BAE RAD750, Vorago VA10820) operate in the $0.3$--$1$\,W range at $100$--$200$\,MHz~\cite{denby}, implying roughly $0.5$\,nJ per cycle. 6FLAS spends about
$132$\,nJ to verify a packet, versus $7.3\,\mu$J for IPsec ESP and $393$\,nJ for
the SRv6 HMAC baseline.

The $20$-bit Flow Label gives $2^{20}\approx1.05$M values per security domain; with $256$ domains
the effective identifier space is $\approx2.7\times10^8$, a factor of $270$ above the
$10^6$ concurrent-flow target. Within one domain, deriving $\fl$ via a keyed PRF spreads flows uniformly, so collisions begin to appear as the
active-flow count approaches $\sqrt{2^{20}}\approx1024$ per domain. At the
maximum per-domain concurrent flow count of $\lfloor10^6/256\rfloor\approx3906$
flows, the expected number of colliding Flow-Label pairs within one domain is
$3906^2/(2\cdot2^{20})\approx7.3$, so at least one collision is likely.
A Flow Label collision is not a security failure because the per-packet MAC of
\cref{eq:tag} still binds the routing context including source and destination addresses; the consequence is that colliding flows share the same flow identifier but remain distinguishable by their addresses. Since the replay window is indexed by $(\domid, \text{direction})$ and operates on the per-domain sequence counter, Flow Label collisions do not affect replay detection (which tracks sequence numbers, not flow labels). Larger deployments can reduce Flow Label collision probability by partitioning traffic across additional domains.
